\documentclass[11pt]{article}
\usepackage[utf8]{inputenc}
\begin{document}
\title{Selección geográfica y generación de ROI de análisis}
\maketitle

\section{Objetivo}
Buscar ubicaciones con Nominatim y confirmar una región de interés (ROI) antes de consultar NISAR GCOV. Las geometrías administrativas originales que caben en los límites configurados se conservan; para las regiones demasiado grandes se construye un rectángulo interior, sin cruzar sus bordes ni huecos.

\section{Uso}
Al escribir al menos tres caracteres en ``Buscar ubicación'', aparecen sugerencias de municipios, ciudades o países con una breve espera entre pulsaciones. Se puede elegir con clic o con las flechas y Enter; Escape cierra la lista. También se puede buscar manualmente con el botón ``Buscar'', incluso si las sugerencias no están disponibles. La lista de resultados muestra nombres de lugares sin clasificaciones técnicas; internamente se retiene la jerarquía para distinguir entidades y se prioriza el límite administrativo frente al punto correspondiente. Al seleccionar el lugar se previsualiza una ROI con radio predeterminado de 5 km (configurable entre 1 y 15 km). Se puede confirmar en un mapa editable o dibujar manualmente. Cuando solo existe un punto, se avisa que la ROI rectangular es aproximada. Si no puede construirse un rectángulo válido dentro del límite administrativo, se informa que debe dibujarse manualmente; no se convierte silenciosamente en una ROI puntual.

En el ``Catálogo GCOV · ASF'', los selectores de fecha ``Desde'' y ``Hasta'' conservan los valores elegidos al cambiar de ROI, mapa o plegar el panel; ambos se pueden limpiar. Si ``Desde'' supera ``Hasta'', aparece un error y se impide la búsqueda. La consulta ASF sigue recibiendo las fechas seleccionadas y la ROI efectiva sin alterar el procesamiento científico.

\section{Implementación}
\texttt{backend/scenes/location\_suggest.py} consulta Photon para sugerir únicamente nombres de lugares; limita la respuesta, filtra tipos, usa caché y una cuota de nuevas consultas compartida por quienes usen la misma caché Django; no fabrica geometrías. \texttt{web/src/components/LocationSearch.tsx} espera 550 ms después de escribir y cancela solicitudes anteriores. Elegir una sugerencia inicia la búsqueda explícita existente; la generación de ROI nunca utiliza una geometría de Photon. Si el servicio de sugerencias falla, el botón de búsqueda sigue disponible.

\texttt{backend/scenes/geosearch.py} solicita \texttt{polygon\_geojson=1}, valida el GeoJSON de Nominatim y obtiene la geometría seleccionada por \texttt{place\_id} desde la búsqueda en el servidor. Si un \texttt{Polygon} o \texttt{MultiPolygon} administrativo cabe en la ventana geodésica local y en 0.6 grados por eje, devuelve la geometría GeoJSON original, con todos sus vértices y huecos. Si no cabe, prueba el punto interior representativo y otros centros de una rejilla determinista: para cada uno reduce el radio hasta que \texttt{covers} de Shapely comprueba la contención del rectángulo completo. Selecciona el rectángulo de mayor radio entre los centros examinados; el radio aplicado puede ser menor que el solicitado. Es una búsqueda acotada de candidatos, no un cálculo global del máximo rectángulo inscrito.

\texttt{web/src/App.tsx} guarda \texttt{effectiveRoi} y metadatos en el mapa editable; \texttt{web/src/hooks/useSavedMaps.ts} conserva la geometría completa al recargar. \texttt{web/src/components/MapCanvas.tsx} usa la misma geometría completa en su fuente de ROI confirmada. La máscara inversa visual representa todas las islas y huecos, además de sus contornos; el anillo \texttt{croppedRegion} sigue disponible para herramientas de dibujo heredadas, pero no sustituye a la ROI efectiva. \texttt{web/src/api/client.ts} envía la ROI efectiva a los puntos de entrada de GCOV y trabajos de ROI, donde \texttt{pipeline/roi\_explore.py} valida \texttt{Polygon} y \texttt{MultiPolygon} con sus anillos originales.

\section{Limitaciones}
El servicio público Photon no garantiza disponibilidad ni límites de uso estables; su falta no bloquea la búsqueda explícita. El Nominatim público no se consulta durante la escritura (su política prohíbe autocompletar). El proveedor de fechas nativo del navegador determina la apariencia del calendario desplegable; no se dispone aquí de comprobación visual interactiva para todos los navegadores. La rejilla de centros para rectángulos interiores es una heurística determinista: puede no encontrar un rectángulo en regiones extremadamente angostas aunque exista uno; en ese caso se solicita la delimitación manual. Las geometrías polares y las que cruzan el antimeridiano se rechazan antes de consultar ASF. El porcentaje GCOV de huella es estimado; no demuestra cobertura de píxeles válidos ni una causa física.

\section{Verificación}
En \texttt{backend}, ejecutar \texttt{PYTHONPATH=.. python manage.py test scenes.test\_location\_suggest scenes.test\_geosearch scenes.test\_gcov}; en \texttt{web}, ejecutar \texttt{node --test tests/*.test.mjs}, \texttt{npm run build} y el lint focalizado. Las 47 pruebas Django de sugerencias, \texttt{scenes.test\_geosearch} y \texttt{scenes.test\_gcov} pasaron; cubren geometrías originales, polígonos cóncavos, islas, huecos y contención con \texttt{covers}, radio de puntos, límite angular y entrada GCOV. Las 51 pruebas web pasaron; cubren la conservación del GeoJSON entre previsualización, persistencia y cuerpo GCOV, la máscara de multipolígono, deduplicación, orden y limpieza de fechas y filtros ASF. La compilación web y el lint focalizado pasaron. El documento compiló con \texttt{pdflatex}. Aún no se han consultado Nominatim/ASF reales ni se ha verificado visualmente el calendario en navegador.
\end{document}
